\noindent\textbf{预条件子正定性证明。}\label{app:invertible}

\textbf{正定性证明：} 不失一般性，假设最大坐标的系统矩阵 $\mathbf{H}$ 是对称正定的。（否则，可以添加 $\alpha \I$ 来确保正定性。）那么其子块 $\mathbf{H}_{\mathbf{p}}$ 和 $\mathbf{H}_{\mathbf{q}}$ 也是对称正定的。

该预条件子基于公式~\ref{eq:preconditioner_structure} 具有以下块结构：
\begin{equation}
  \begin{aligned}
    \mathbf{P} & =\begin{bmatrix}
                    \sum_{i=0}^{n}\mathbf{h}_{i} & \mathbf{0}                                                                                                                                                \\
                    \mathbf{0}                   & \overbrace{\operatorname{blk-diag}(\J_{\p,\omega}^\top \mathbf{H}_\p \J_{\p,\omega})+ \J_{\q,\omega}^\top \mathbf{H}_\q \J_{\q,\omega}}^{\mathbf{P}_{22}}
                  \end{bmatrix}.
  \end{aligned}
\end{equation}
其中 $\mathbf{h}_{i}$ 是大于 $\mathbf{0}$ 的对角矩阵。

显然，第一个对角块是正定的，因此 $\mathbf{P}$ 正定当且仅当 $\mathbf{P}_{22}$ 正定。

根据公式~\ref{eq:A}，可以指出 $\J_{\mathbf{p, s}}^{\top}\mathbf{H}_{\mathbf{p}}\J_{\mathbf{p, s}}$ 和 $\J_{\mathbf{q, s}}^{\top}\mathbf{H}_{\mathbf{q}}\J_{\mathbf{q, s}}$ 至少都是半正定的，因为：
\begin{equation}
  \forall \mathbf{v} \ne \mathbf{0}, \quad
  \mathbf{v}^{\top}\J_{*}^{\top}\mathbf{H}_{*}\underbrace{\J_{*} \mathbf{v}}_{\mathbf{y}} = \mathbf{y}^{\top}
  \mathbf{H}_{*}\mathbf{y} \geq 0.
\end{equation}
因此，$\operatorname{blk-diag}\left(\J_{\p, \omega}^\top \mathbf{H}_{\p} \J_{\p,
    \omega}\right)$ 仍然是半正定的。

接下来证明 $\J_{\mathbf{q,\omega}}^{\top}\mathbf{H}_{\mathbf{q}}\J_{\mathbf{q,\omega}}$ 是正定的。

从 $\omega$ 到 $\mathbf{q}$ 的映射 $\mathcal{T}_{\mathbf{q,\omega}}$ 是局部单射映射，因此其雅可比 $\J_{\mathbf{q,\omega}}$ 具有满列秩。考虑到 $\mathbf{H_{\mathbf q}}$ 的正定性，可将其分解为 $\mathbf{LL}^{\top}$，并且：
\begin{equation}
  \begin{aligned}
    \mathrm{rank}(\J_{\mathbf{q,\omega}}^{\top}\mathbf{H}_{\mathbf{q}}\J_{\mathbf{q,\omega}})
     & =
    \mathrm{rank}(\J_{\mathbf{q,\omega}}^{\top}\mathbf{LL}^{\top}\J_{\mathbf{q,\omega}}) \\
     & =\mathrm{rank}(\mathbf{L}^{\top}\J_{\mathbf{q,\omega}})                           \\
     & =\mathrm{rank}(\J_{\mathbf{q,\omega}})
     & =\mathrm{dim}(\omega).
  \end{aligned}
\end{equation}

因此，$\J_{\mathbf{q,\omega}}^{\top}\mathbf{H}_{\mathbf{q}}\J_{\mathbf{q,\omega}}$ 仍然是正定的。

至此，该预条件子是正定的：
\begin{equation}
  \begin{aligned}
    \mathbf{P} & =\left(\operatorname{blk-diag}\left(\J_{\p, \s}^\top
    \mathbf{H}_{\p} \J_{\p, \s}\right) + \begin{bmatrix} \mathbf{0} &
    \mathbf{0}    \\ \mathbf{0} & \J_{\q,\omega}^\top \mathbf{H}_\q \J_{\q,
      \omega}\end{bmatrix}\right)                                                                                                                                                                                                                                                  \\
               & =\begin{bmatrix}
                    \overbrace{\sum_{i=0}^{n}\mathbf{h}_{i}}^{\text{正定}} & \mathbf{0}                                                                                                                                                                                                                                               \\
                    \mathbf{0}                                                          & \underbrace{\underbrace{\operatorname{blk-diag}(\J_{\p,\omega}^\top \mathbf{H}_\p \J_{\p,\omega})}_{\text{半正定}}+ \underbrace{\J_{\q,\omega}^\top \mathbf{H}_\q \J_{\q,\omega}}_{\text{正定}}}_{\text{正定}}
                  \end{bmatrix}.
  \end{aligned}
\end{equation}
